% DP30 SHA256 32eed2415e3ac16810d3e3e13973131df45e83e948979ba7395f6f698ff49b3a
\begin{proposition}[DP30: a dual direction with a unit $p$-th derivative]\label{DP30}
Let $A$ be a complete mixed-characteristic discrete valuation ring with perfect residue field $k$. Let $G/A$ be a one-dimensional strict formal $\mathcal O$-module, where $\mathcal O\subset A$ is the ring of integers of the unramified degree-$h$ extension of $\Q_p$, and suppose its special fibre has height exactly $h\ge2$. Let $H=G[p^\infty]^D$ be its actual Cartier dual and $\widehat H$ its connected formal group. Then $\Lie(H_k)$ has a basis $e_0,\ldots,e_{h-2}$ with
\[
 e_i^{[p]}=e_{i+1}\ (i<h-2),\qquad e_{h-2}^{[p]}=0.
\]
Thus the restricted $p$-operation has kernel dimension one and image dimension $h-2$. If $h>2$, there are a primitive integral invariant derivation $D$ on $\widehat H$ and an integral formal coordinate $f$ such that $D^pf(0)\in A^\times$.
\end{proposition}
\begin{proof}
Teichm\"uller averaging as in DP19 gives an integral coordinate $x$ with $[\zeta](x)=\zeta x$ for every $\zeta^{p^h-1}=1$. The group law then satisfies $F(\zeta X,\zeta Y)=\zeta F(X,Y)$, so its nonzero monomials have total degree congruent to one modulo $q-1$, where $q=p^h$. In particular $F=X+Y+$ terms of degree at least $q$, integrally and after Frobenius twisting.

For the next calculation work over $k$. Relative Frobenius $F_G:G\to G^{(p)}$ and Verschiebung $V_G:G^{(p)}\to G$ compose in either order to multiplication by $p$; their kernel orders are $p$ and $p^{h-1}$. These facts and the duality of $F,V$ are Tate \cite[\S2.2, diagram (3), pp.~163--164; \S2.3, Proposition 3, p.~166]{Tate67}. Restrict $V_G$ to $G^{(p)}[p]$. Its image has order $p$ and lies in $G[F_G]$, so it is exactly $G[F_G]$. Exact finite Cartier duality \cite[\S1.2--1.3, pp.~159--160]{Tate67} gives
\[
 H[F_H]^D\simeq G[p]/G[F_G].
\]
Height $h$ implies $[p](x)=u(x)x^q$ with $u(0)\ne0$, so $G[p]=G[F_G^h]$ as subgroup schemes, by equality of their defining ideals. If $r=p^{h-1}$, relative Frobenius induces
\[
 G[p]/G[F_G]\simeq G^{(p)}[F^{h-1}].
\]
On coordinates this is $k[z]/(z^r)\hookrightarrow k[x]/(x^q)$, $z\mapsto x^p$. The target is free of rank $p$ with basis $1,x,\ldots,x^{p-1}$, proving faithful flatness and the asserted quotient. In $k[X,Y]/(X^r,Y^r)$ every surviving monomial has degree at most $2r-2<q$. The nonlinear part of the twisted group law vanishes. Hence
\begin{equation}\label{eq:dual-F-kernel}
 H[F_H]^D\simeq\alpha_{p^{h-1}}.
\end{equation}

Put $P=H[F_H]$. Its distribution Hopf algebra is, by duality,
\[
 \Dist(P)=k[t]/(t^r),\qquad\Delta(t)=t\otimes1+1\otimes t.
\]
The Lie algebra is its primitive subspace: the augmentation Leibniz rule for a tangent functional is precisely the primitive-element identity in the dual Hopf algebra. This subspace has basis $t,t^p,\ldots,t^{p^{h-2}}$. Indeed these powers are primitive, while for any other exponent $1\le n<r$ a nonzero base-$p$ digit gives an interior binomial coefficient $\binom n{p^j}\not\equiv0\pmod p$. The resulting mixed tensor term survives; different total degrees cannot cancel it. Constants are not primitive.

The restricted operation is the $p$-th power in this algebra: convolution powers of a tangent functional correspond to powers of its invariant derivation. It shifts the displayed basis and kills its last member. Since $dF_H=0$, the inclusion $P\hookrightarrow H$ identifies their Lie algebras and their restricted operations. Perfectness of $k$ gives the stated image and kernel dimensions.

The connected formal group exists by Tate \cite[\S2.2, Proposition 1, p.~162 and (4), p.~164]{Tate67}; its tangent lattice is free of rank $h-1$, since dual dimensions sum to height by Proposition 3. For $h>2$, lift a tangent vector $v$ with $v^{[p]}\ne0$ to this lattice and let $D$ be its invariant derivation. The lift is primitive. Some integral coordinate $f$ has differential detecting $v^{[p]}$. Reduction commutes with $D$ and its iterates, so
\[
 \overline{D^pf(0)}=df_0(v^{[p]})\ne0.
\]
Thus $D^pf(0)$ is a unit. No assertion that the full $H$ is connected is needed for this last argument.
\end{proof}

% DP31 SHA256 84bf7014febae0988bb54b2055d90cce452cef0e6cd97a376ddd630ddd6d81d0
\begin{theorem}[DP31: the scalar higher-height correction is integral]\label{DP31}
Let $h>2$, let $k$ contain $\F_q$, and now take $R=W(k)$ and $K=R[1/p]$. Under (H2)--(H4),
\begin{equation}\label{eq:scalar-gauge}
 B(S,T):=\exp\beta(L_h(S),L_h(T))\in1+ST\,R[[S,T]].
\end{equation}
Consequently the scalar function of DP29 satisfies the stronger estimate
\begin{equation}\label{eq:scalar-full-slope}
 v_p(\kappa![\mathbf t^\kappa]\Theta)
 \ge c_h|\kappa|-\frac{\sum_i s_p(\kappa_i)}{p-1}.
\end{equation}
The slope $c_h$ is optimal within this class, by \eqref{eq:sharp-corrected-scalar}. This theorem concerns two scalar copies of the height-$h$ group, not the genuinely dual pair.
\end{theorem}
\begin{proof}
Let $G$ have law $F_h$, and write $R_n=R/p^nR$. We use the established formal biextension theorem: Zink \cite[\S4.4, Proposition 121 and (243), pp.~241, 246]{Zink} gives an injection
\[
 \operatorname{Biext}^1(G_n\times G_n,\widehat{\Gm})
 \hookrightarrow\Hom(G_n,G_n^t).
\]
The formal groups here come from displays; the existence/equivalence applies over the Artinian local rings $R_n$ by \cite[Theorem 9, or Theorem 103]{Zink}. The dual display gives the Cartier-dual formal group. Both $G$ and its Cartier dual are connected for $h>1$, with special-fibre slopes $1/h$ and $(h-1)/h$ respectively. The hypotheses of the injection are therefore satisfied; we do not need the additional surjectivity in Zink's Theorem 122.

First, the integral units $C_X=H_X$ and $C_Y=H_Y$ define a biextension without presupposing \eqref{eq:scalar-gauge}. Over $K$ they are the two multiplicative coboundaries of $B$. Therefore they satisfy symmetry, associativity in each variable, and the interchange identity. These identities hold over $R$, since $R\hookrightarrow K$. On the trivial formal multiplicative torsor over $G\times G$, define the first partial product by
\[
 (S,T,u)\star_X(S',T,u')
  =(F_h(S,S'),T,uu'C_X(S,S',T)),
\]
and define the second analogously with $C_Y$. The identities just stated make a biextension. A normalized splitting is exactly a unit $I$ with
\begin{equation}\label{eq:gauge-split}
 \frac{I(F_h(S,S'),T)}{I(S,T)I(S',T)}=C_X,\qquad
 \frac{I(S,F_h(T,T'))}{I(S,T)I(S,T')}=C_Y,
 \qquad I(S,0)=I(0,T)=1.
\end{equation}

Second, $\Hom(G_n,G_n^t)=0$. Over $k$, distinct slopes give zero rational Hom, by the semisimple isocrystal classification after faithful extension to an algebraic closure \cite[\S1.3, first paragraph of p.~120]{Katz79}, hence zero integral Hom. The integral Hom embeds in its rationalization: a homomorphism killed by $p^a$ is zero since multiplication by $p^a$ on a $p$-divisible group is surjective as a sheaf. To lift the vanishing, let $f:G_n\to G_n^t$ reduce to zero, and suppose its formal-coordinate components are in $p^jR_n[[T]]$, $1\le j<n$. Write $f=p^jg$ componentwise. The target $p$-series has linear part $p$ times the identity and higher terms of order at least two, so
\[
 f([p]_G(T))=[p]_{G^t}(f(T))\equiv0\pmod {p^{j+1}}.
\]
After division by $p^j$ and reduction, each component satisfies $\bar g([p]_{G_k}(T))=0$. Substitution of the nonzero series $[p]_{G_k}(T)$, whose initial degree is $q$, is injective on $k[[T]]$. Thus $f$ is divisible by $p^{j+1}$, and induction proves $f=0$.

The biextension injection now yields a splitting $I_n\in1+ST\,R_n[[S,T]]$ at every $n$. We verify compatibility rather than assume that independently chosen splittings lift. Two normalized splittings differ by a bihomomorphism $J$, which obeys
\[
 J([p]_G(S),T)=J(S,T)^p.
\]
Modulo $p$, let $b>0$ be the smallest $T$-degree in $J-1$, if it is nonzero. Its coefficient on the left remains nonzero under substitution of $[p]_{G_k}(S)$; on the right every positive $T$-degree is at least $pb>b$. This contradiction proves $J=1$ modulo $p$. If $J-1=p^jV$ for $j\ge1$, the right side minus one is divisible by $p^{j+1}$. Dividing and reducing gives $\bar V([p]_{G_k}(S),T)=0$, and injectivity of this substitution on $k[[S,T]]$ proves the next divisibility. Hence $J=1$ over every $R_n$.

Thus the normalized splittings are unique and automatically compatible. Their coefficientwise inverse limit is $I\in1+ST\,R[[S,T]]$ satisfying \eqref{eq:gauge-split}. In additive coordinates,
\[
 A(X,Y)=\beta(X,Y)-\log I(E_h(X),E_h(Y))
\]
is separately additive, by equality of the two coboundaries. In characteristic zero, interior binomial coefficients force a separately additive series to be $aXY$. The logarithm and exponential are Teichm\"uller-equivariant. Consequently the assumed torus symmetry of $\beta$ makes $I_n(\zeta S,\zeta^{-p}T)$ another normalized splitting. Uniqueness makes $I$, and hence $A$, invariant. But $XY$ has character $\zeta^{1-p}\ne1$, so $a=0$. This proves $B=I$ and \eqref{eq:scalar-gauge}.

Finally, write the correction factor in $\Theta$ as
\[
 \prod_j B(E_h(\ell_j),E_h(\epsilon\ell_j))^{\sigma/2}.
\]
Each binomial power is an integral series because $p$ is odd. The whole expression is an integral function of finitely many inverse-logarithm coordinates of weight at least $\rho_h$. Therefore $v_p([\mathbf t^\kappa]\Theta)\ge-\rho_h|\kappa|$. Legendre's factorial formula gives \eqref{eq:scalar-full-slope}. The explicit nonzero example in DP29 attains the slope along $q$-power orders.
\end{proof}

% DP32 SHA256 8a4b04104fe06d196f3ec0402270a7326c7931b3ded661bc28ded6a8702075a4
\begin{theorem}[DP32: sharp asymptotic slopes on the actual Honda dual]\label{DP32}
Let $p$ be any prime, $h\ge2$, $q=p^h$, $d=h-1$, and $\mathcal O=W(\F_q)$. The Cartier dual $H=G_h^D$ of the Honda group with logarithm $L_h(T)=\sum_{n\ge0}T^{q^n}/p^n$ has integral coordinates $X_1,\ldots,X_d$ and normalized primitive invariant derivations $D_i$ for which, setting
\[
 c_i=\frac{p^i}{q-1},\qquad r_i=\frac1{p-1}-c_i\quad(1\le i\le d),
\]
every integral formal function satisfies coefficientwise
\begin{equation}\label{eq:dual-mixed}
 v_p(D_1^{a_1}\cdots D_d^{a_d}f)
 \ge\max\left\{0,\ \sum_i c_i a_i-\frac{\sum_i s_p(a_i)}{p-1}\right\}.
\end{equation}
For every $i$ and $n=q^r$,
\begin{equation}\label{eq:dual-sharp}
 v_p(D_i^nX_i(0))=c_i(n-1).
\end{equation}
Thus every displayed directional slope is optimal, and $\sum_i c_i=c_h$. These assertions persist over $W(k)\supset\mathcal O$. In particular the order-$p$ obstruction DP30 does not imply absence of positive asymptotic gain on the dual.
\end{theorem}
\begin{proof}
We construct an integral group and prove its dual identification before estimating it. Over $\Z_p$, define a vector of monomials
\[
 M_1(\mathbf X)=X_d^{p^2},\qquad M_i(\mathbf X)=X_{i-1}^p\ (i\ge2),
 \qquad \mathcal L(\mathbf X)=\sum_{n\ge0}\frac{M^{\circ n}(\mathbf X)}{p^n}.
\]
For $d=1$, this means $M_1=X_1^{p^2}$. The degrees tend to infinity, so the series is well defined and has identity linear term. Its functional equations are
\[
 \mathcal L_1=X_1+p^{-1}\mathcal L_d(\mathbf X^{p^2}),\qquad
 \mathcal L_i=X_i+p^{-1}\mathcal L_{i-1}(\mathbf X^p)\ (i\ge2).
\]
Apply Hazewinkel \cite[Lemma 2.3, (2.2)--(2.6), p.~74]{Hazewinkel} with $B=\Z_p$, ideal $pB$, $\sigma=\id$, the lemma's $q=p$, matrix $s_1$ having $1/p$ just below the diagonal, and $s_2$ having $1/p$ in entry $(1,d)$. When $d=1$, $s_1=0$. It gives an integral commutative law $F_H=\mathcal L^{-1}(\mathcal L(\mathbf X)+\mathcal L(\mathbf Y))$ and an integral $p$-series. The exact identity
\[
 \mathcal L(M(\mathbf X))=p(\mathcal L(\mathbf X)-\mathbf X)
\]
and the congruence clause \cite[(2.8), p.~75]{Hazewinkel} imply $[p]_H\equiv M\pmod p$.

Here the finite-flat condition can be checked directly. Put $C=\Z_p[[U_1,\ldots,U_d]]$, $B=\Z_p[[X_1,\ldots,X_d]]$ and send $U_i$ to $[p]_{H,i}(\mathbf X)$. Modulo $p$, $B/p$ is free over $C/p$ on the $p^h$ monomials
\[
 X_1^{b_1}\cdots X_d^{b_d},\qquad
 0\le b_i<p\ (i<d),\quad 0\le b_d<p^2.
\]
Their lifts define $C^{p^h}\to B$. Successively lifting the coefficient expansion modulo $p$ and dividing the remainder by $p$ proves surjectivity by $p$-adic completeness. A relation has all coefficients divisible by $p$; dividing by $p$ is valid since $B$ is $p$-torsion-free. Repetition and $p$-adic separatedness prove injectivity. Thus $[p]_H$ is finite free of degree $p^h$. Its iterated kernels form a connected $p$-divisible group of height $h$ and dimension $d$, as in Tate \cite[\S2.2, Proposition 1, pp.~162--164]{Tate67}.

Everything is defined over $\Z_p$. On the special fibre over $\F_p$, a full cycle of $M$ has total Frobenius exponent $2+(d-1)=h$, so $[p]_H^d=M^d=F_H^h$. For $G=H^D$, dual dimensions give $\dim G=1$ \cite[\S2.3, Proposition 3, p.~166]{Tate67}, and duality gives $V_G^h=[p]_G^d$. Since $F_G^hV_G^h=[p]_G^h$ and $d=h-1$, cancellation gives $F_G^h=[p]_G$. Cancellation is valid because multiplication by $p^d$ is surjective as an fppf sheaf. All these maps are over $\F_p$, so they are actual endomorphisms with no unidentified twists. Their Frobenius kernels show that $G$ is connected. It therefore has an integral one-dimensional formal group over $\Z_p$ by Tate's equivalence. In an integral coordinate $T$ its $p$-series $f$ has linear term $pT$ and reduction $T^q$.

Over $\mathcal O=W(\F_q)$, this is precisely a Lubin--Tate series for uniformizer $p$. The known uniqueness and strict-isomorphism theorem \cite[\S5, pp.~76--77, (5.1)--(5.4)]{Hazewinkel} identifies the group with this $p$-series with the normalized Lubin--Tate group $F_f$, and identifies all such groups for the same uniformizer integrally and strictly. The Honda group $G_h$ has a $p$-series in this class. Thus $G\simeq G_h$ integrally over $\mathcal O$, and dualizing identifies the constructed $H$ with the actual Cartier dual claimed in the statement. No guessed dual logarithm or unproved lifting of an action is involved.

We now compute its radius in each coordinate. Put $S_h=(q-1)/(p-1)$, choose $\pi^{q-1}=p$, and set $\alpha_i=\pi^{S_h-p^i}$, so $v_p(\alpha_i)=r_i>0$. Scale coordinate $i$ in input and output by $\alpha_i$. For every arrow $X_j^{p^e}$ in $M_i$ one has $p^er_j-r_i=1$. Consequently the $n=1$ term of
\[
 \mathcal L'_i(\mathbf Z)=\alpha_i^{-1}\mathcal L_i(\alpha_1Z_1,\ldots,\alpha_dZ_d)
\]
has coefficient exactly one. If an $n$-fold cyclic iterate has successive exponents $e_1,\ldots,e_n$, its coefficient valuation after scaling is
\[
 1+p^{e_1}+p^{e_1+e_2}+\cdots+p^{e_1+\cdots+e_{n-1}}-n.
\]
This is positive for $n\ge2$. Hence $\mathcal L'$ is integral and reduces to $\mathbf Z+M(\mathbf Z)$. Its inverse $\mathcal E'$ is integral. In characteristic $p$, the monomial map $M$ is additive, and therefore
\[
 \overline{\mathcal E'}(\mathbf Z)=\sum_{n\ge0}(-1)^nM^{\circ n}(\mathbf Z).
\]
Composition with $\id+M$ verifies the cancellation, also at $p=2$. In component $i$, after setting other variables to zero, terms survive exactly at cyclic step numbers $dr$, with degrees $q^r$ and nonzero coefficients $(-1)^{dr}$. Undoing the input and output scales gives
\[
 v_p([t^{q^r}]\mathcal E_i(t e_i))=-(q^r-1)r_i.
\]

The normalized invariant derivations are integral and commute. In additive logarithm coordinates they are partial derivatives, so
\[
 D^{\mathbf a}f(0)=\mathbf a![\mathbf t^{\mathbf a}]f(\mathcal E(\mathbf t)).
\]
Integrality of $\mathcal E'$ says that every coefficient of $\mathcal E_j$ has weighted valuation at least $r_j$ when variable $t_i$ is given weight $r_i$. Integral substitution therefore bounds $[\mathbf t^{\mathbf a}]f(\mathcal E)$ by $-\sum_i r_i a_i$. Adding Legendre's exact factorial valuation proves \eqref{eq:dual-mixed} at zero; integrality also supplies the bound zero. Apply the same computation in $\mathbf Y$ to the integral translated function $f(F_H(\mathbf X,\mathbf Y))$, treating $\mathbf X$ as coefficient variables. Translation invariance gives the coefficientwise bound. At $n=q^r$, the exact exponential coefficient above and $s_p(n)=1$ prove \eqref{eq:dual-sharp}. Finally
\[
 \sum_{i=1}^{h-1}c_i=\frac{p+\cdots+p^{h-1}}{p^h-1}=c_h.
\]
This proves all assertions, including sharpness against a larger slope with sublinear error.
\end{proof}

% DP33 SHA256 79168ae11d5d8e2af3acd39c1b56e3820ed50aaca6c52eb1e9fc3a3320fe6bc4
\begin{theorem}[DP33: optimal affine errors for the height-three dual]\label{DP33}
In the two-coordinate model of DP32 with $h=3$, write $X=X_1$, $Y=X_2$, $q=p^3$, $c_X=p/(q-1)$ and $c_Y=p^2/(q-1)$. For every integral $f$ the coefficientwise bounds improve to
\[
 v_p(D_X^nf)\ge\max\{0,\ceil{c_Xn-1}\},\qquad
 v_p(D_Y^nf)\ge\max\{0,\ceil{c_Yn-1}\}.
\]
Both slopes and both constant errors $1$ are optimal. Mixed derivatives have the sum of the two displayed integer exponents as a lower bound. In addition to the exact diagonal sequences in DP32,
\[
 v_p(D_X^{pq^r}Y(0))=\frac{p^2(q^r-1)}{q-1},\qquad
 v_p(D_Y^{p^2q^r}X(0))=\frac{p(q^{r+1}-1)}{q-1}.
\]
These statements include $p=2$.
\end{theorem}
\begin{proof}
The actual-dual model and its integral coordinates are already established by DP32. Its logarithm begins
\[
 \mathcal L_1=X+Y^{p^2}/p+X^q/p^2+\cdots,\qquad
 \mathcal L_2=Y+X^p/p+Y^q/p^2+\cdots.
\]
Let $\Delta_{a,b}$ be the integral coefficient operators of translation by $(S,T)$. They commute with one another and with $D_X,D_Y$, by commutativity of the group. In characteristic zero the translation operator is $\exp(\mathcal L_1(S,T)D_X+\mathcal L_2(S,T)D_Y)$. Comparing coefficients gives
\[
 \Delta_{p,0}=D_X^p/p!+D_Y/p,\qquad
 \Delta_{0,p^2}=D_Y^{p^2}/(p^2)!+D_X/p,\qquad
 \Delta_{0,p}=D_Y^p/p!.
\]
Consequently, with integral commuting operators $T,S$ and units $a,b$,
\begin{equation}\label{eq:dual-low-ops}
 D_X^p=aD_Y+pT,\quad a=-(p-1)!,\qquad
 D_Y^{p^2}=p^p bD_X+p^{p+1}S,\quad b=-\frac{(p^2)!}{p^{p+1}}.
\end{equation}
Also $D_Y^p$ is divisible by $p$ as an operator. Raising the first identity to $p$ shows $D_X^{p^2}$ is divisible by $p$: the leading term uses $D_Y^p$, and every remaining binomial term has a factor $p$. Thus for all $n$ the elementary bounds are $\floor{n/p^2}$ in direction $X$ and $\floor{n/p}$ in direction $Y$.

Raise the first identity in \eqref{eq:dual-low-ops} to $p^2$. Its nonleading term indexed by $1\le j\le p^2$ has valuation at least
\[
 v_p\binom{p^2}{j}+j+\floor{(p^2-j)/p}\ge p+1.
\]
For $j=p^2$ this is $p^2\ge p+1$. Otherwise $v_p\binom{p^2}{j}=2-v_p(j)$. If $j=ap+b$ with $1\le b<p$, the expression is $p+1+a(p-1)+b$; if $j=ap$, $1\le a<p$, it is $p+1+a(p-1)$. Applying the second identity to the leading term proves
\[
 D_X^q=p^p u_XD_X+p^{p+1}U_X
\]
with $u_X$ a unit and $U_X$ integral. Raise the second identity of \eqref{eq:dual-low-ops} to $p$. Its leading term uses the first identity; every mixed term $1\le j<p$ has valuation at least $1+p(p-j)+(p+1)j=p^2+j+1$, and its last term has valuation $p(p+1)\ge p^2+1$. Hence
\[
 D_Y^q=p^{p^2}u_YD_Y+p^{p^2+1}U_Y.
\]

For either direction put $(r,B,c)=(p^2,p,c_X)$ or $(p,p^2,c_Y)$. Then $c=B/(q-1)<1$ and $cr=q/(q-1)$. For $0\le n<q$, the maximum of $cn-\floor{n/r}$ is one, attained at $n=q-1$: within a block it increases, and successive block maxima increase by $cr-1>0$. This proves the affine inequality at the initial orders. For $n\ge q$, apply the common recurrence $D^q=p^BuD+p^{B+1}U$ to $D^{n-q}f$. Induction bounds the first term by $cn-1$ and the second by $cn-1+(1-c)$. Integrality and integer valuations give the ceiling. Since the bounds are on the whole formal function ring, composing the two operator containments proves the mixed assertion.

The affine constants are attained by the same monomial $f=X^{p-1}Y^{p^2-1}$ at $n=q-1$. Along the $X$ logarithmic axis the leading terms of the exponential are $X=t+\cdots$, $Y=-t^p/p+\cdots$, so $f$ has first degree $q-1$ with coefficient $(-1/p)^{p^2-1}$. Since $v_p((q-1)!)=p^2+p-2$, the derivative has valuation $p-1=c_X(q-1)-1$. Along the $Y$ axis the leading terms are $X=-t^{p^2}/p+\cdots$, $Y=t+\cdots$; the coefficient is $(-1/p)^{p-1}$ and the derivative valuation is $p^2-1=c_Y(q-1)-1$.

For the additional exact sequences, DP32's scaled inverse here reduces to
\[
 \overline{\mathcal E'}_1(u,v)=\sum_{r\ge0}u^{q^r}-\sum_{r\ge0}v^{p^2q^r},\qquad
 \overline{\mathcal E'}_2(u,v)=\sum_{r\ge0}v^{q^r}-\sum_{r\ge0}u^{pq^r}.
\]
The nonzero off-diagonal coefficients, the respective input/output weights $r_X=(p^2+1)/(q-1)$ and $r_Y=(p+1)/(q-1)$, and Legendre's formula give exactly the two stated valuations. The diagonal sequences from DP32 prove slope optimality. In particular $D_X^pY(0)=-(p-1)!$ is a unit, while $c_X>0$.
\end{proof}

% DP34 SHA256 eb5be984ea9f34075f00480ee997f6d19af47421ba58b48ed0feb9328502c110
\begin{theorem}[DP34: optimal constant-one affine bounds on every Honda dual]\label{DP34}
In the actual-dual coordinates of DP32, for every $h\ge2$, every $i\in\{1,\ldots,h-1\}$ and every integral $f$,
\[
 D_i^nf\in p^{b_i(n)}R[[\mathbf X]],\qquad
 b_i(n)=\max\{0,\ceil{c_i n-1}\},\qquad c_i=\frac{p^i}{p^h-1}.
\]
Both the slope and the constant $1$ are optimal. With $q=p^h$ and $d=h-1$, the same monomial
\[
 f_*(\mathbf X)=\left(\prod_{j<d}X_j^{p-1}\right)X_d^{p^2-1}
\]
satisfies $v_p(D_i^{q-1}f_*(0))=p^i-1$ for every $i$. Mixed derivatives have the sum $\sum_i b_i(a_i)$ as a lower bound. The displayed ideal containments persist under integral base change.
\end{theorem}
\begin{proof}
Work first over $\Z_p$ in the model already identified by DP32. Integral translation coefficients and the first nonlinear terms of its logarithm give, exactly as in DP33,
\begin{equation}\label{eq:all-dual-low}
 D_i^p=aD_{i+1}+pT_i\quad(i<d),\qquad
 D_d^{p^2}=p^pbD_1+p^{p+1}T_d,\qquad D_d^p\in p\operatorname{End}_{\rm int}.
\end{equation}
Here $a=-(p-1)!$, $b=-(p^2)!/p^{p+1}$ are units; all remainder operators preserve the integral function ring and commute. The last two identities also hold when $d=1$. Descending induction, raising the first identity to $p^{h-i-1}$, shows that $D_i^{R_i}\in p\operatorname{End}_{\rm int}$ for $R_i=p^{h-i}$. Hence $D_i^n$ has the preliminary divisibility $\floor{n/R_i}$.

Fix $i<d$, put $B=p^i$, and start with $D_i^q$. Successively substitute the first identity of \eqref{eq:all-dual-low} to move its leading term through directions $j=i,i+1,\ldots,d-1$. At the step $j=i+t$, this identity is raised to $N=p^{h-t-1}$, and $N/R_{j+1}=B$. Its $m$-th nonleading term has valuation at least
\[
 v_p\binom Nm+m+\floor{(N-m)/R_{j+1}}
 =B+v_p\binom Nm+m-\ceil{m/R_{j+1}}.
\]
For $1\le m<N$, the binomial coefficient is divisible by $p$ because $N$ is a $p$-power, and $m-\ceil{m/R_{j+1}}\ge0$. For $m=N$ the bound is $N\ge p^{i+1}\ge B+1$. Thus every discarded term belongs to $p^{B+1}\operatorname{End}_{\rm int}$. The leading term at the end is a unit times $D_d^{p^{i+1}}$. Raise the second identity of \eqref{eq:all-dual-low} to $p^{i-1}$ and factor out $p^B$. Modulo $p$, what remains is a unit times $D_1^{p^{i-1}}$, which the first identities modulo $p$ identify with a unit times $D_i$. Therefore
\begin{equation}\label{eq:all-dual-q}
 D_i^q=p^{p^i}u_iD_i+p^{p^i+1}U_i
\end{equation}
with a unit $u_i$ and an integral operator $U_i$. For $i=d$, use the same last step directly by raising the second identity to $p^{d-1}$. This includes $d=1$.

Now $c_iR_i=q/(q-1)$ and $c_i(q-1)=p^i$. For $0\le n<q$, the increasing block maxima give $\floor{n/R_i}\ge c_i n-1$, with equality in the real difference at $n=q-1$. For $n\ge q$, applying \eqref{eq:all-dual-q} to $D_i^{n-q}f$ bounds its first term by $c_i n-1$ and its second by $c_i n-1+(1-c_i)\ge c_i n-1$. This proves the affine bound by induction. All coefficients and operators were over $\Z_p$, so integer valuations give the asserted ideal containment there. Base change preserves that containment; no fractional valuation is rounded on a ramified base. Composition proves the mixed bound.

It remains to prove the sharp common witness, including exclusion of possible lower-degree terms. The explicit logarithm and its inverse satisfy the Teichm\"uller equivariance
\[
 \mathcal L_j(\zeta^pX_1,\zeta^{p^2}X_2,\ldots,\zeta^{p^d}X_d)
 =\zeta^{p^j}\mathcal L_j(\mathbf X),\qquad\zeta^{q-1}=1.
\]
For the ordinary arrows this is $p\cdot p^{j-1}=p^j$; for the exceptional arrow it is $p^2p^d\equiv p\pmod{q-1}$. Thus a nonzero degree $n$ in $\mathcal E_j(te_i)$ requires $np^i\equiv p^j\pmod{q-1}$. Its smallest possible positive degree is
\[
 m_{ij}=p^{j-i}\ (j\ge i),\qquad m_{ij}=p^{h+j-i}\ (j<i).
\]
DP32's scaled inverse gives a nonzero coefficient at this first path degree, with exact valuation $r_j-m_{ij}r_i$. Equivariance excludes all smaller degrees. If $e_j=p-1$ for $j<d$ and $e_d=p^2-1$, then
\[
 \sum_j e_jm_{ij}=q-1,\qquad \sum_j e_jr_j=d.
\]
The first equality telescopes: the ranges $i\le j<d$, $j=d$, and $j<i$ contribute $p^{d-i}-1$, $p^{h-i+1}-p^{d-i}$, and $q-p^{h-i+1}$ respectively. For the second, use $\sum e_j=(p-1)(h+p-1)$ and $\sum e_jp^j=p(q-1)$. The unique lowest-degree term of $f_*(\mathcal E(te_i))$ therefore has degree $q-1$ and valuation $d-(q-1)r_i$. Finally
\[
 v_p((q-1)!)=\frac{q-1}{p-1}-h
\]
makes the derivative valuation $c_i(q-1)-h+d=p^i-1$. This proves the optimal constant in every direction. Slope optimality is already the exact sequence in DP32. Every step applies at $p=2$ and $h=2$.
\end{proof}

% DP36 SHA256 6d271f98b2f952c441994eb8379baa439ec05e21399683769713962379bb93c0
\begin{proposition}[DP36: unequal weights restore the kernel quantization]\label{DP36}
Let $p$ be odd, $h\ge2$, $q=p^h$, $d=h-1$, and $R=W(k)$ with $\F_q\subset k$. Use the actual pair $G_h,H=G_h^D$ of DP32, with logarithms $L_0,\mathcal L$ and inverses $E_0,\mathcal E$. Put
\[
 s=(p-1)^{-1},\quad\rho=(q-1)^{-1},\quad r_0=\rho,\quad r_i=s-p^i\rho\ (1\le i\le d).
\]
Suppose $\beta(x,\mathbf y)\in R[1/p][[x,\mathbf y]]$ vanishes on either axis, has no bidegree-$(1,1)$ monomial, and satisfies
\[
 \beta(\zeta x,\zeta^{-p}y_1,\ldots,\zeta^{-p^d}y_d)=\beta(x,\mathbf y)
\]
for a primitive $(q-1)$-st root $\zeta$. Suppose its two exponential addition coboundaries
\[
 C_X=\exp\bigl(\delta_x\beta(L_0(S),L_0(S'),\mathcal L(\mathbf T))\bigr),\quad
 C_Y=\exp\bigl(\delta_{\mathbf y}\beta(L_0(S),\mathcal L(\mathbf T),\mathcal L(\mathbf T'))\bigr)
\]
are integral units with constant coefficient one. Here $\delta_{\mathbf y}$ uses addition of the whole vector in logarithm coordinates. For any integral series $U$ in the group coordinates and $\sigma\in\{1,-1\}$, set
\[
 \Theta(x,\mathbf y)=U(E_0(x),\mathcal E(\mathbf y))\exp\bigl((\sigma/2)\beta(x,\mathbf y)\bigr).
\]
Then
\begin{equation}\label{eq:true-pair-weight}
 v_p(\mathbf a![x^{a_0}\mathbf y^{\mathbf a'}]\Theta)
 \ge (s-p\rho)a_0+\sum_{i=1}^d\left(\frac{p^{i+1}}{q-1}-1\right)a_i
       -\frac{\sum_{i=0}^d s_p(a_i)}{p-1}.
\end{equation}
On the specific path $x=t$, $y_d=\epsilon t$, $y_i=0$ for $i<d$, with $\epsilon\in R^\times$, this gives
\begin{equation}\label{eq:true-pair-path}
 v_p(n![t^n]\Theta(t,0,\ldots,0,\epsilon t))
 \ge\rho n-\frac{s_p(n)}{p-1}.
\end{equation}
These are sufficient estimates, with no asserted optimality or geometric realization of the hypotheses or path.
\end{proposition}
\begin{proof}
If $x^a\mathbf y^{\mathbf b}$ is invariant, $a-\sum_i p^ib_i$ is divisible by $q-1$. Its weighted degree is
\[
 r_0a+\sum_i r_ib_i
 =\rho\left(a-\sum_i p^ib_i\right)+s|\mathbf b|\in s\Z.
\]
All $r_i$ are positive. Every nonconstant invariant monomial therefore has weighted degree at least $s$. The same calculation applies with duplicated $x$-variables or duplicated $\mathbf y$-vectors. Integral coefficient valuations are nonnegative, so the lower bound remains valid. The actual logarithms, exponentials and group laws are equivariant for these characters by their explicit formulas in DP32. Thus the integral series $C_X-1,C_Y-1$ have weighted valuation at least $s$.

The inverse-logarithm estimates give weight at least $r_0$ to $E_0$ and at least $r_i$ to $\mathcal E_i$; substituting them preserves this bound. Taking the formal logarithm consequently gives
\[
 w_{\mathbf r}(\delta_x\beta),\ w_{\mathbf r}(\delta_{\mathbf y}\beta)\ge s.
\]
Indeed its $m$-th term has weight at least $ms-v_p(m)\ge s$, since $v_p(m)\le(m-1)/(p-1)$.

For a monomial $x^a\mathbf y^{\mathbf b}$ of $\beta$, both side degrees are positive and not both one. If $a\ge2$, choose an interior binomial coefficient of valuation at most one as in DP29 and read the corresponding coefficient of $\delta_x\beta$. If $a=1$ and $|\mathbf b|\ge2$, either two components of $\mathbf b$ are nonzero, in which case a split sending one whole component to the first vector has coefficient one, or a single component has exponent at least two, in which case use the same binomial choice in that component. Recombining the split exponents recovers the original monomial uniquely, so there is no contamination by another coefficient. It follows that $w_{\mathbf r}(\beta)\ge s-1$.

Each monomial of $\beta$ has positive weighted degree at least $s$. Increasing all weights from $r_i$ to $pr_i$ therefore adds at least $(p-1)s=1$ to every coefficient weight. Thus $w_{p\mathbf r}(\beta)\ge s$, and its exponential with multiplier $\sigma/2$ has nonnegative weight by $v_p(m!)\le ms$. The integral $U$ factor has nonnegative weight as well; increasing the positive weights preserves the established inverse bounds. Hence the ordinary coefficient is at least $-p\sum_i r_ia_i$ in valuation. Adding Legendre's exact factorial identity proves \eqref{eq:true-pair-weight}.

For the stated path assign $t$ weight $pr_d$. Since $r_d\ge r_0$, both nonzero substitutions have at least their required weights. Setting the other coordinates to zero preserves the bound. The resulting factorial coefficient is bounded by $(s-pr_d)n-s_p(n)/(p-1)$. The identity $s-pr_d=\rho$ proves \eqref{eq:true-pair-path}.
\end{proof}
